\section{Motivation and Challenges}
\label{sec:motivation}

Efficient vector search is critical for large-scale systems used in recommendation, semantic search, and information retrieval. These applications demand the ability to process a high volume of nearest neighbor queries with low latency, even as the underlying data evolves. To meet these requirements, vector databases—such as Milvus \cite{wang2021milvus}, Pinecone \cite{pinecone-url}, AnalyticDB-V \cite{wei2020analyticdb}, VBASE \cite{zhang2023vbase}, and Qdrant \cite{noauthor_qdrant_nodate}—utilize specialized vector indexes (e.g., Faiss-IVF, HNSW, Vamana) that support fast approximate nearest neighbor (ANN) queries. However, serving these dynamic workloads introduces significant challenges in maintaining query performance and accuracy as data and query patterns shift over time.


\subsection{Vector Search Workload}
A vector search workload is a continuous, evolving stream of \textbf{queries} and \textbf{updates}:
\begin{itemize}[leftmargin=*, itemsep=0pt, topsep=0pt]
    \item \textbf{Queries:} Given a query vector $q$, the goal is to find the top-$k$ nearest neighbors in a set $\mathbf{X}$. Exact linear search is too slow for large, high-dimensional datasets, so ANN indexes are used. These indexes approximate nearest neighbors with controlled recall to lower latency by orders of magnitude.
    \item \textbf{Updates:} The dataset evolves over time. Insertions add new vectors representing fresh content (e.g., new products, trending news articles), and deletions remove outdated entries. Typically, updates are applied in a batched fashion.
\end{itemize}

\textbf{Recall@k} is the standard metric for accuracy, defined as:
$\frac{|\mathbf{G} \cap \mathbf{R}|}{k}$ where \(\mathbf{R}\) is the vectors returned by the approximate search, and \(\mathbf{G}\) is the ground truth set. Maintaining a consistent recall target (e.g., $>90\%$) and low latency (e.g., milliseconds per query) as both data and query patterns shift is a key challenge. The complexity of these workloads stems from their inherently dynamic and skewed nature, which few existing indexing methods handle gracefully.

\subsection{Why Real-World Workloads are Hard}

\paragraph{Skewed Read Patterns}
In practice, user queries concentrate on popular items. For example, queries against a Wikipedia-derived dataset tend to focus on a small subset of entities at any given time. As a result, certain partitions or graph regions of the index are accessed disproportionately often.

\paragraph{Skewed Write Patterns}
Insertions and deletions are also rarely uniform. New data often arrives in bursts—e.g., new Wikipedia pages added monthly, new products introduced ahead of a shopping season, or newly relevant embeddings generated by continuously updated language models. 

\begin{figure}[t]
    \centering
    \begin{subfigure}[t]{0.48\linewidth}
        \centering
        \includegraphics[width=\linewidth]{sections/figures/wiki_distributions.pdf}
        \caption{Read (top) and write skew.}
        \label{fig:wiki_combined}
    \end{subfigure}%
    \hfill
    \begin{subfigure}[t]{0.48\linewidth}
        \centering
        \includegraphics[width=\linewidth]{sections/figures/latency_and_recall_per_operation_ivf.pdf}
        \caption{Query performance.}
        \label{fig:ivf_latency}
    \end{subfigure}
    \caption{Skewed access patterns of Faiss-IVF index partitions in the \textsc{Wikipedia-12M} workload and their effect on query performance for Faiss-IVF and SCANN}
    \label{fig:combined_figures}
\end{figure}



\paragraph{Real-World Example: Wikipedia-12M} In our evaluation, we prepared \textsc{Wikipedia-12M}, a workload based on a subset of Wikipedia articles derived from publicly available monthly pageview statistics \cite{wiki}. Over 103 months, the dataset grows from millions to tens of millions of vectors. Popular articles dominate query traffic, while embeddings of newly created pages accumulate in certain regions of the embedding space. This workload shows read skew and write skew, as evidenced by Figure \ref{fig:wiki_combined}, reads and writes predominantly affect a small portion of the index. 

\begin{table}[t]
\small
\centering
\caption{\edit{Comparison of updatable vector indexes.} \textbf{Tuning}: Requires manual parameter tuning in indexing/query processing. \textbf{Maintenance}: Modifies index with incremental updates. \textbf{Adaptive}: Uses query information to inform indexing.}
\label{tab:method_comparison}
\begin{tabular}{lccc}
\toprule
\textbf{Method} & \textbf{Tuning} & \textbf{Maint.} & \textbf{Adaptive} \\
\midrule
Quake \textbf{(Ours)} & \xmark & \cmark & \cmark \\
Faiss-IVF~\cite{douze2024faisslibrary} & \cmark & \xmark & \xmark \\
DeDrift~\cite{baranchuk_dedrift_2023} & \cmark & \cmark & \xmark \\
SpFresh~\cite{xu_spfresh_2023} & \cmark & \cmark & \xmark \\
SCANN~\cite{guo_accelerating_2020, soar_2023} & \cmark & \cmark & \xmark \\
DiskANN~\cite{ni_diskann_2023, singh2021freshdiskann} & \cmark & \cmark & \xmark \\
Faiss-HNSW~\cite{10.1109/TPAMI.2018.2889473} & \cmark & \xmark & \xmark \\
SVS~\cite{aguerrebere2024locally} & \cmark & \cmark & \xmark \\
\bottomrule
\end{tabular}
\end{table}

\subsection{Shortcomings of Existing Approaches}
Existing indexes were often developed and evaluated under assumptions of static data distributions; conditions not met in real-world use cases. Table~\ref{tab:method_comparison} compares a range of state-of-the-art vector indexes. Although widely adopted in vector databases, none fully solve the problem of maintaining low-latency, high-recall search under dynamic, skewed workloads without constant manual intervention or offline tuning.

\paragraph{Graph Indexes}  
Graph-based index systems, such as HNSW~\cite{10.1109/TPAMI.2018.2889473}, DiskANN~\cite{ni_diskann_2023, singh2021freshdiskann}, and SVS ~\cite{aguerrebere2024locally} construct a proximity graph where each node represents a vector connected to its approximate neighbors. These indexes achieve high recall with low latency in static settings by efficiently traversing the graph to locate nearest neighbors using a process known as \emph{greedy traversal}. However, maintaining the graph structure under frequent updates is computationally intensive, as each update may require rewiring multiple edges to preserve graph connectivity and proximity properties. Our evaluation (Table \ref{tab:method_comparison}) shows that update latency can be multiple orders of magnitude higher than partitioned indexes.

\paragraph{Partitioned Indexes}  
Partitioned indexes such as Faiss-IVF~\cite{douze2024faisslibrary}, SCANN~\cite{guo_accelerating_2020, soar_2023}, and SpFresh~\cite{xu_spfresh_2023} divide the vector space into disjoint partitions using a clustering algorithm such as k-means. Queries are processed by scanning a subset of partitions to retrieve approximate nearest neighbors. Partitioned indexes are more update-friendly than graph-based method since insertions and deletions leads to sequential access. For write skewed workloads some partitions become significantly larger, degrading query latency, this can be exacerbated by read skew if large partitions are also more frequently accessed ("hot partitions"). Query processing is \emph{memory-bound}, as achieving high recall requires scanning many megabytes of data across multiple partitions, for example reaching a recall target of 90\% on the \textsc{MSTuring100M} dataset requires each query to scan 1GB of vectors. Moreover, most partitioned indexes use a fixed number of partitions to probe (\emph{nprobe}), which does not adapt to changing data distributions or query patterns, leading to either insufficient recall or excessive data scanning. The challenges yield subpar performance for partitioned indexes on real-world workloads. For example, Figure \ref{fig:ivf_latency} shows the degradation of latency and recall over time when using Faiss-IVF and SCANN with a fixed nprobe on \textsc{Wikipedia-12M} (workload details in Section~\ref{sec:experiments}).

\edit{\paragraph{Early Termination}}
Early-termination methods have been proposed to reduce query latency or meet recall targets in partitioned indexes by dynamically adjusting the number of partitions scanned per query. \textbf{SPANN}~\cite{chen_spann_nodate} applies a simple rule: it prunes partitions once the centroid distance exceeds a user-tuned threshold relative to the closest centroid. \textbf{LAET}~\cite{li2020improving} is a learning-based approach that predicts the required \emph{nprobe} per query using a trained model, but still requires dataset-specific training and calibration for each recall target. \textbf{Auncel}~\cite{zhang2023fast} uses a geometric model to estimate when recall for a given query, setting \emph{nprobe} per query, but its conservative estimation leads to substantial overshooting of the recall target (See Figure 13 in \cite{zhang2023fast}). All three methods require tuning or calibration and do not adapt to changes in the index structure or data distribution.

\subsection{Technical Challenges for Partitioned Indexes}
The following technical challenges are yet to be solved by existing partitioned indexes
\begin{enumerate}[leftmargin=*, itemsep=-.25em, topsep=0pt]
    \item \textbf{Adaptation to Queries} Query adaptivity is overlooked by existing partitioned index approaches and exhibits an opportunity for optimization, particularly for maintaining hot partitions induced by read skew.
    \item \textbf{Online Adjustment of Nprobe} As the index structure and data change, partitioned indexes need to adjust the number of partitions scanned or recall will suffer. \edit{Existing early termination works are insufficient as they assume a static index and require retuning as the index and data change.}
    \item \textbf{Performance Gap with Graph Indexes} Standard partitioned indexes such as Faiss-IVF are memory bound, and exhibit an order of magnitude higher query latency in comparison than graph indexes.
\end{enumerate}

Quake is our solution to these technical challenges. Quake A) adapts the index structure to queries by utilizing maintenance that minimizes a cost model for query latency, B) using a recall estimation model, Quake individually sets nprobe for queries to meet recall targets as the index structure changes, and C) uses NUMA-aware parallelism in order to saturate memory bandwidth during query processing, closing the performance gap with graph indexes. We next briefly discuss other related work before covering Quake in detail.
